\MajorSection{القسم السابع}

\Couplet{١٦٣}{ليس للإنسان روح؛ إنها حال من الأحوال، لا تزال لا شيء، صوتًا وكلمة؛}{تُولّد على هذا النحو شيئًا ذا جوهر، حتى ترى العين ما سمعته الأذن.}{}

\Couplet{١٦٤}{أين كانت روحه حين كان الوحش الهمجي الذي يهيم في الغابات الأولى؛}{وأي هيئة كانت لها، وأي مقام، وأي دور أدّته في خطة الطبيعة؟}{}

\Couplet{١٦٥}{هذه الروح صُنعت لحل لغز؛ ومن يريد هذه الثنائية العبثية؛}{ألا تكفيني ذاتي؟ وما حاجتي إلى «أنا» داخل «أنا»؟}{}

\Couplet{١٦٦}{كلمات، كلمات تولّد أشياء! فالروح وافدة جديدة على المشهد؛}{ألا يكفي نفَس الحياة لتشغيل الآلة المولودة من المادة؟}{}

\Couplet{١٦٧}{نعرف نشأة الروح، ونتتبّعها إلى ساعة ميلادها؛}{ونرصد نموها كما نما البشر، حتى تباهوا بأنهم السادة الوحيدون للأرض.}{}

\Couplet{١٦٨}{جرى شوط الوجود منذ فجر الحياة في مسار لا انقطاع فيه؛}{وما يطيب للناس أن يسمّوه أرواحهم بدأ في الخنزير والكلب.}{}

\Couplet{١٦٩}{الحياة سلّم لا تنتهي درجاته، يخفي درجاته عن أعين الإنسان؛}{تقوم قدمه في ظلام العماء، ويرتفع رأسه عاليًا فوق السماوات.}{}

\Couplet{١٧٠}{لا تحمل سلسلة الوجود أي انقطاع؛ كل الأشياء بدأت في وحدة؛}{وتنتظم الحلقات في صف، وإن لم ير أحد، ربما، تعاقبها.}{}

\Couplet{١٧١}{الشبح، ذلك الخوف الطبيعي المتجسد من الموت الكئيب والتعفن القذر؛}{ولّد الروح والنفس والظل، مع حشد هاديس الشاحب الواهن.}{}

\Couplet{١٧٢}{واحتاجت النفس إلى نفس أعظم، نفس للنفوس تحكم الجمع؛}{ومن هنا جاءت القوى الروحية ومراتبها، وكلها ولدها الشبح الهمجي.}{}

\Couplet{١٧٣}{ليست لكم، يا أهل الكتاب، هذه الرؤى الخرافية الحسنة المحببة؛}{ولدتها آلهة أرض خيمي، وسافرت بعيدًا وراء البحار!}{\BurtonNote{حاشية برتون: مصر؛ كام وكِم وخِم في الهيروغليفية، وخيمي في الديموطيقية.}}

\Couplet{١٧٤}{«ذهن الإنسان الفاني الخالد!» نسمع المتعصب ذا الرئتين الصاخبتين يصيح؛}{وذهنه لا يعني إلا مجموع فكره، وخلاصة «أنا» ذرية.}{}

\Couplet{١٧٥}{الفكر عمل الدماغ والعصب؛ فقير وضئيل في الأبله الصغير الجمجمة؛}{مريض في المرض، نائم في النوم، وميت حين يُسدل الموت ستار المشهد.}{}

\Couplet{١٧٦}{«هُراء!» يقول الزاهد: «نعرف جيدًا تعليم المدرسة الممقوتة؛}{التي تجعل الإنسان آلة ذاتية الحركة، والذهن إفرازًا، والروح كلمة».}{}

\Couplet{١٧٧}{«سرعان ما تثرثرون، يا تجّار المادة، عن الجزيئات والبروتوبلازم؛}{وعن تطور بقعة هلام، وقرود ارتقت إلى مرتبة الإنسان».}{}

\Couplet{١٧٨}{اعتراض عبثي! كل ما هو كائن جاء إما بمعجزة وإما بناموس؛}{فلماذا تهدرون على هذا كراهيتكم وخوفكم، وعلى ذاك حبكم وخشوعكم؟}{}

\Couplet{١٧٩}{لماذا تكدّسون على كلمة كل هذه الكراهية، ولماذا تنسبون إلى المثال نموذجًا أوليًا؛}{ولماذا تحمّلون المادة روحًا؟ أتحتاج خطتكم إلى ملحق؟}{}

\Couplet{١٨٠}{أليس أعلى الشرف لمن استخرج الأفضل من الأسوأ؛}{أما يجوز لصانعكم أن يصنع العالم من مادة، إن وافق ذلك أمره؟}{}

\Couplet{١٨١}{بل كلما كانت المادة أوضع، كانت يد الصانع أمهر؛}{فكفّوا عن تقييد قدرتكم الكلية نفسها، وتحديدها، وفهمها.}{}

\Couplet{١٨٢}{«العقل والغريزة!» ما أشد حبنا للعب بكلمات ترضي كبرياءنا؛}{ونحاول أن نخفي بها، بألقاب كاذبة مصطنعة، الأصل الوضيع لجنسنا النبيل!}{}

\Couplet{١٨٣}{وأطلب منكم أن تقرؤوا بدل «العطية الإلهية»: العمل الأفضل للدماغ الأعلى؛}{الذي يختلف عن الغريزة في الدرجة، كما يختلف منجم الذهب عن عرق الرصاص.}{}

\Couplet{١٨٤}{العقل هو حَكَم الحياة الوحيد، وخيط المتاهة السحرية الوحيد؛}{وثمة عوالم فوقه، وراء علمه؛ ولكن ما يناقضه لا يمكن أبدًا أن يكون حقًا.}{}

\Couplet{١٨٥}{«يندفع الحمقى حيث تخشى الملائكة أن تخطو!» وللملائكة والحمقى الحق نفسه}{في أن يفعلوا ما تأمرهم به الطبيعة، بلا رجاء مدح، ولا خوف ذم!}{}
